\IfFileExists{stacks-project.cls}{%
\documentclass{stacks-project}
}{%
\documentclass{amsart}
}

% For dealing with references we use the comment environment
\usepackage{verbatim}
\newenvironment{reference}{\comment}{\endcomment}
%\newenvironment{reference}{}{}
\newenvironment{slogan}{\comment}{\endcomment}
\newenvironment{history}{\comment}{\endcomment}

% For commutative diagrams we use Xy-pic
\usepackage[all]{xy}

% We use 2cell for 2-commutative diagrams.
\xyoption{2cell}
\UseAllTwocells

% We use multicol for the list of chapters between chapters
\usepackage{multicol}

% This is generally recommended for better output
\usepackage{lmodern}
\usepackage[T1]{fontenc}

% For cross-file-references
\usepackage{xr-hyper}

% Package for hypertext links:
\usepackage{hyperref}

% For any local file, say "hello.tex" you want to link to please
% use \externaldocument[hello-]{hello}
\externaldocument[introduction-]{introduction}
\externaldocument[conventions-]{conventions}
\externaldocument[sets-]{sets}
\externaldocument[categories-]{categories}
\externaldocument[topology-]{topology}
\externaldocument[sheaves-]{sheaves}
\externaldocument[sites-]{sites}
\externaldocument[stacks-]{stacks}
\externaldocument[fields-]{fields}
\externaldocument[algebra-]{algebra}
\externaldocument[brauer-]{brauer}
\externaldocument[homology-]{homology}
\externaldocument[derived-]{derived}
\externaldocument[simplicial-]{simplicial}
\externaldocument[more-algebra-]{more-algebra}
\externaldocument[smoothing-]{smoothing}
\externaldocument[modules-]{modules}
\externaldocument[sites-modules-]{sites-modules}
\externaldocument[injectives-]{injectives}
\externaldocument[cohomology-]{cohomology}
\externaldocument[sites-cohomology-]{sites-cohomology}
\externaldocument[dga-]{dga}
\externaldocument[dpa-]{dpa}
\externaldocument[sdga-]{sdga}
\externaldocument[hypercovering-]{hypercovering}
\externaldocument[schemes-]{schemes}
\externaldocument[constructions-]{constructions}
\externaldocument[properties-]{properties}
\externaldocument[morphisms-]{morphisms}
\externaldocument[coherent-]{coherent}
\externaldocument[divisors-]{divisors}
\externaldocument[limits-]{limits}
\externaldocument[varieties-]{varieties}
\externaldocument[topologies-]{topologies}
\externaldocument[descent-]{descent}
\externaldocument[perfect-]{perfect}
\externaldocument[more-morphisms-]{more-morphisms}
\externaldocument[flat-]{flat}
\externaldocument[groupoids-]{groupoids}
\externaldocument[more-groupoids-]{more-groupoids}
\externaldocument[etale-]{etale}
\externaldocument[chow-]{chow}
\externaldocument[intersection-]{intersection}
\externaldocument[pic-]{pic}
\externaldocument[weil-]{weil}
\externaldocument[adequate-]{adequate}
\externaldocument[dualizing-]{dualizing}
\externaldocument[duality-]{duality}
\externaldocument[discriminant-]{discriminant}
\externaldocument[derham-]{derham}
\externaldocument[local-cohomology-]{local-cohomology}
\externaldocument[algebraization-]{algebraization}
\externaldocument[curves-]{curves}
\externaldocument[resolve-]{resolve}
\externaldocument[models-]{models}
\externaldocument[functors-]{functors}
\externaldocument[equiv-]{equiv}
\externaldocument[pione-]{pione}
\externaldocument[etale-cohomology-]{etale-cohomology}
\externaldocument[proetale-]{proetale}
\externaldocument[relative-cycles-]{relative-cycles}
\externaldocument[more-etale-]{more-etale}
\externaldocument[trace-]{trace}
\externaldocument[crystalline-]{crystalline}
\externaldocument[spaces-]{spaces}
\externaldocument[spaces-properties-]{spaces-properties}
\externaldocument[spaces-morphisms-]{spaces-morphisms}
\externaldocument[decent-spaces-]{decent-spaces}
\externaldocument[spaces-cohomology-]{spaces-cohomology}
\externaldocument[spaces-limits-]{spaces-limits}
\externaldocument[spaces-divisors-]{spaces-divisors}
\externaldocument[spaces-over-fields-]{spaces-over-fields}
\externaldocument[spaces-topologies-]{spaces-topologies}
\externaldocument[spaces-descent-]{spaces-descent}
\externaldocument[spaces-perfect-]{spaces-perfect}
\externaldocument[spaces-more-morphisms-]{spaces-more-morphisms}
\externaldocument[spaces-flat-]{spaces-flat}
\externaldocument[spaces-groupoids-]{spaces-groupoids}
\externaldocument[spaces-more-groupoids-]{spaces-more-groupoids}
\externaldocument[bootstrap-]{bootstrap}
\externaldocument[spaces-pushouts-]{spaces-pushouts}
\externaldocument[spaces-chow-]{spaces-chow}
\externaldocument[groupoids-quotients-]{groupoids-quotients}
\externaldocument[spaces-more-cohomology-]{spaces-more-cohomology}
\externaldocument[spaces-simplicial-]{spaces-simplicial}
\externaldocument[spaces-duality-]{spaces-duality}
\externaldocument[formal-spaces-]{formal-spaces}
\externaldocument[restricted-]{restricted}
\externaldocument[spaces-resolve-]{spaces-resolve}
\externaldocument[formal-defos-]{formal-defos}
\externaldocument[defos-]{defos}
\externaldocument[cotangent-]{cotangent}
\externaldocument[examples-defos-]{examples-defos}
\externaldocument[algebraic-]{algebraic}
\externaldocument[examples-stacks-]{examples-stacks}
\externaldocument[stacks-sheaves-]{stacks-sheaves}
\externaldocument[criteria-]{criteria}
\externaldocument[artin-]{artin}
\externaldocument[quot-]{quot}
\externaldocument[stacks-properties-]{stacks-properties}
\externaldocument[stacks-morphisms-]{stacks-morphisms}
\externaldocument[stacks-limits-]{stacks-limits}
\externaldocument[stacks-cohomology-]{stacks-cohomology}
\externaldocument[stacks-perfect-]{stacks-perfect}
\externaldocument[stacks-introduction-]{stacks-introduction}
\externaldocument[stacks-more-morphisms-]{stacks-more-morphisms}
\externaldocument[stacks-geometry-]{stacks-geometry}
\externaldocument[moduli-]{moduli}
\externaldocument[moduli-curves-]{moduli-curves}
\externaldocument[examples-]{examples}
\externaldocument[exercises-]{exercises}
\externaldocument[guide-]{guide}
\externaldocument[desirables-]{desirables}
\externaldocument[coding-]{coding}
\externaldocument[obsolete-]{obsolete}
\externaldocument[fdl-]{fdl}
\externaldocument[index-]{index}

% Theorem environments.
%
\theoremstyle{plain}
\newtheorem{theorem}[subsection]{Theorem}
\newtheorem{proposition}[subsection]{Proposition}
\newtheorem{lemma}[subsection]{Lemma}

\theoremstyle{definition}
\newtheorem{definition}[subsection]{Definition}
\newtheorem{example}[subsection]{Example}
\newtheorem{exercise}[subsection]{Exercise}
\newtheorem{situation}[subsection]{Situation}

\theoremstyle{remark}
\newtheorem{remark}[subsection]{Remark}
\newtheorem{remarks}[subsection]{Remarks}

\numberwithin{equation}{subsection}

% Macros
%
\def\lim{\mathop{\mathrm{lim}}\nolimits}
\def\colim{\mathop{\mathrm{colim}}\nolimits}
\def\Spec{\mathop{\mathrm{Spec}}}
\def\Hom{\mathop{\mathrm{Hom}}\nolimits}
\def\Ext{\mathop{\mathrm{Ext}}\nolimits}
\def\SheafHom{\mathop{\mathcal{H}\!\mathit{om}}\nolimits}
\def\SheafExt{\mathop{\mathcal{E}\!\mathit{xt}}\nolimits}
\def\Sch{\mathit{Sch}}
\def\Mor{\mathop{\mathrm{Mor}}\nolimits}
\def\Ob{\mathop{\mathrm{Ob}}\nolimits}
\def\Sh{\mathop{\mathit{Sh}}\nolimits}
\def\NL{\mathop{N\!L}\nolimits}
\def\CH{\mathop{\mathrm{CH}}\nolimits}
\def\proetale{{pro\text{-}\acute{e}tale}}
\def\etale{{\acute{e}tale}}
\def\QCoh{\mathit{QCoh}}
\def\Ker{\mathop{\mathrm{Ker}}}
\def\Im{\mathop{\mathrm{Im}}}
\def\Coker{\mathop{\mathrm{Coker}}}
\def\Coim{\mathop{\mathrm{Coim}}}

% Boxtimes
%
\DeclareMathSymbol{\boxtimes}{\mathbin}{AMSa}{"02}

%
% Macros for moduli stacks/spaces
%
\def\QCohstack{\mathcal{QC}\!\mathit{oh}}
\def\Cohstack{\mathcal{C}\!\mathit{oh}}
\def\Spacesstack{\mathcal{S}\!\mathit{paces}}
\def\Quotfunctor{\mathrm{Quot}}
\def\Hilbfunctor{\mathrm{Hilb}}
\def\Curvesstack{\mathcal{C}\!\mathit{urves}}
\def\Polarizedstack{\mathcal{P}\!\mathit{olarized}}
\def\Complexesstack{\mathcal{C}\!\mathit{omplexes}}
% \Pic is the operator that assigns to X its picard group, usage \Pic(X)
% \Picardstack_{X/B} denotes the Picard stack of X over B
% \Picardfunctor_{X/B} denotes the Picard functor of X over B
\def\Pic{\mathop{\mathrm{Pic}}\nolimits}
\def\Picardstack{\mathcal{P}\!\mathit{ic}}
\def\Picardfunctor{\mathrm{Pic}}
\def\Deformationcategory{\mathcal{D}\!\mathit{ef}}
\usepackage{cleveref}
\crefname{originalreference}{original reference}{original references}
\Crefname{originalreference}{Original reference}{Original references}
\makeatletter
\newlabel{algebra-lemma-dimension-base-fibre-equals-total}{{00ON}{source}{}{sourcepointer-1}{}}
\newlabel{algebra-lemma-dimension-base-fibre-equals-total@cref}{{[originalreference][0][]00ON}{[1][0][]source}{}{}{}}
\makeatother
\crefname{originalreference}{original reference}{original references}
\Crefname{originalreference}{Original reference}{Original references}
\makeatletter
\newlabel{algebra-lemma-faithfully-flat-universally-injective}{{05CK}{source}{}{sourcepointer-2}{}}
\newlabel{algebra-lemma-faithfully-flat-universally-injective@cref}{{[originalreference][0][]05CK}{[1][0][]source}{}{}{}}
\makeatother
\crefname{originalreference}{original reference}{original references}
\Crefname{originalreference}{Original reference}{Original references}
\makeatletter
\newlabel{algebra-lemma-regular-ring-CM}{{00NQ}{source}{}{sourcepointer-3}{}}
\newlabel{algebra-lemma-regular-ring-CM@cref}{{[originalreference][0][]00NQ}{[1][0][]source}{}{}{}}
\makeatother
\crefname{originalreference}{original reference}{original references}
\Crefname{originalreference}{Original reference}{Original references}
\makeatletter
\newlabel{algebra-lemma-regular-sequence-in-polynomial-ring}{{07DW}{source}{}{sourcepointer-4}{}}
\newlabel{algebra-lemma-regular-sequence-in-polynomial-ring@cref}{{[originalreference][0][]07DW}{[1][0][]source}{}{}{}}
\makeatother
\crefname{originalreference}{original reference}{original references}
\Crefname{originalreference}{Original reference}{Original references}
\makeatletter
\newlabel{algebra-lemma-regular-sequence-powers}{{07DV}{source}{}{sourcepointer-5}{}}
\newlabel{algebra-lemma-regular-sequence-powers@cref}{{[originalreference][0][]07DV}{[1][0][]source}{}{}{}}
\makeatother
\crefname{originallemma}{original lemma}{original lemmas}
\Crefname{originallemma}{Original lemma}{Original lemmas}
\makeatletter
\newlabel{lemma-compare-standard}{{07EZ}{source}{}{sourcepointer-6}{}}
\newlabel{lemma-compare-standard@cref}{{[originallemma][0][]07EZ}{[1][0][]source}{}{}{}}
\makeatother
\crefname{originallemma}{original lemma}{original lemmas}
\Crefname{originallemma}{Original lemma}{Original lemmas}
\makeatletter
\newlabel{lemma-delocalize-height-zero}{{07FA}{source}{}{sourcepointer-7}{}}
\newlabel{lemma-delocalize-height-zero@cref}{{[originallemma][0][]07FA}{[1][0][]source}{}{}{}}
\makeatother
\crefname{originallemma}{original lemma}{original lemmas}
\Crefname{originallemma}{Original lemma}{Original lemmas}
\makeatletter
\newlabel{lemma-enlarge-solution-modulo}{{07FI}{source}{}{sourcepointer-8}{}}
\newlabel{lemma-enlarge-solution-modulo@cref}{{[originallemma][0][]07FI}{[1][0][]source}{}{}{}}
\makeatother
\crefname{originallemma}{original lemma}{original lemmas}
\Crefname{originallemma}{Original lemma}{Original lemmas}
\makeatletter
\newlabel{lemma-improve-presentation}{{07CE}{source}{}{sourcepointer-9}{}}
\newlabel{lemma-improve-presentation@cref}{{[originallemma][0][]07CE}{[1][0][]source}{}{}{}}
\makeatother
\crefname{originallemma}{original lemma}{original lemmas}
\Crefname{originallemma}{Original lemma}{Original lemmas}
\makeatletter
\newlabel{lemma-lift-solution}{{07F8}{source}{}{sourcepointer-10}{}}
\newlabel{lemma-lift-solution@cref}{{[originallemma][0][]07F8}{[1][0][]source}{}{}{}}
\makeatother
\crefname{originallemma}{original lemma}{original lemmas}
\Crefname{originallemma}{Original lemma}{Original lemmas}
\makeatletter
\newlabel{lemma-ogoma}{{07FC}{source}{}{sourcepointer-11}{}}
\newlabel{lemma-ogoma@cref}{{[originallemma][0][]07FC}{[1][0][]source}{}{}{}}
\makeatother
\crefname{originallemma}{original lemma}{original lemmas}
\Crefname{originallemma}{Original lemma}{Original lemmas}
\makeatletter
\newlabel{lemma-solution-modulo}{{07FH}{source}{}{sourcepointer-12}{}}
\newlabel{lemma-solution-modulo@cref}{{[originallemma][0][]07FH}{[1][0][]source}{}{}{}}
\makeatother
\crefname{originallemma}{original lemma}{original lemmas}
\Crefname{originallemma}{Original lemma}{Original lemmas}
\makeatletter
\newlabel{lemma-strictly-standard-base-change}{{07CC}{source}{}{sourcepointer-13}{}}
\newlabel{lemma-strictly-standard-base-change@cref}{{[originallemma][0][]07CC}{[1][0][]source}{}{}{}}
\makeatother
\crefname{originalreference}{original reference}{original references}
\Crefname{originalreference}{Original reference}{Original references}
\makeatletter
\newlabel{situation-local}{{07F7}{source}{}{sourcepointer-14}{}}
\newlabel{situation-local@cref}{{[originalreference][0][]07F7}{[1][0][]source}{}{}{}}
\makeatother
\title{The general resolution lemma for geometrically regular targets in positive characteristic}
\author{The Stacks Project Authors}
\date{stacks-project:smoothing.tex:local-snapshot-20261012; GNU FDL 1.2 or later}
\begin{document}
\maketitle
\noindent\textit{Editorial note.} The source marks one step '(minor detail omitted)' in Ad (item-choose-deltai), namely the clearing-denominators choice of s realizing both displayed conditions; it is the mechanical denominator step of the construction and is carried as an explicit core obligation.\par\medskip
\noindent\textit{Editorial note.} Inputs lemma-improve-presentation (07CE), lemma-compare-standard (07EZ) and lemma-solution-modulo (07FH) are published entries used as standing inputs by their display numbers; not proven on the author's behalf.\par\medskip
\noindent\textit{Editorial note.} Steps the author leaves to the reader are flagged here for the reviewer: source line(s) 3182 carry an "omitted"/"follows easily"/"immediate from the definitions" phrase. Those sentences are retained verbatim; the displayed argument around them stays inside the counted span.\par\medskip
\noindent\textit{Allowed standard prerequisites (editorial).}\begin{itemize}
\item Presentation improvement with retraction. lemma-improve-presentation (same chapter, PUBLISHED 07CE). Standardizes the auxiliary algebra B and gives the section C -> B.
\item Comparison of standardness conditions. lemma-compare-standard (same chapter, PUBLISHED 07EZ). Promotes smoothness of C\_\{x\_i\} with free differentials to strict standardness of x\_i\textasciicircum{}c.
\item Artinian approximation factorization. lemma-solution-modulo (same chapter, PUBLISHED 07FH). Supplies the step (item-NP) factorization k[y]\_p/p\textasciicircum{}n -> D -> Lambda\_q/q\textasciicircum{}n.
\item Ogoma's annihilator lemma. lemma-ogoma (carded L26814). Used in the steps (item-modify-pii) and (item-choose-deltai) to enforce the annihilator equalities.
\item Enlarging the Artinian approximation ring. lemma-enlarge-solution-modulo (carded L26817). Supplies the multipliers delta\_i of step (item-choose-deltai) with the (a)(b)(c) properties.
\item Removing strictly standard elements from a resolution. lemma-lift-solution (carded L26812). Justifies the reduction (item-second-resolve) modulo the powers gamma\_i\textasciicircum{}e.
\item Height-zero delocalization. lemma-delocalize-height-zero (carded L26813). Justifies the reduction (item-third-resolve) to the localization at r.
\item Strictly standard elements survive base change. lemma-strictly-standard-base-change (carded L26796). Carries strict standardness of gamma\_i\textasciicircum{}c from x\_i\textasciicircum{}c to the base-changed algebra.
\end{itemize}\medskip
\ifdefined\hypertarget\hypertarget{sourcepointer-1}{}\fi
\noindent\textit{Source reference.} Original reference 00ON: focused-seeds/geometry/S9-smo-resolve-general/ref-tags.txt, source lines 5--5. Navigation pointer only: the counted proof cites this statement by \textbackslash{}ref as a standing input; the deduction the counted proof has to make is carried out inside the counted span, and no part of that deduction is delegated to this pointer.\par\medskip
\ifdefined\hypertarget\hypertarget{sourcepointer-2}{}\fi
\noindent\textit{Source reference.} Original reference 05CK: focused-seeds/geometry/S9-smo-resolve-general/ref-tags.txt, source lines 6--6. Navigation pointer only: the counted proof cites this statement by \textbackslash{}ref as a standing input; the deduction the counted proof has to make is carried out inside the counted span, and no part of that deduction is delegated to this pointer.\par\medskip
\ifdefined\hypertarget\hypertarget{sourcepointer-3}{}\fi
\noindent\textit{Source reference.} Original reference 00NQ: focused-seeds/geometry/S9-smo-resolve-general/ref-tags.txt, source lines 7--7. Navigation pointer only: the counted proof cites this statement by \textbackslash{}ref as a standing input; the deduction the counted proof has to make is carried out inside the counted span, and no part of that deduction is delegated to this pointer.\par\medskip
\ifdefined\hypertarget\hypertarget{sourcepointer-4}{}\fi
\noindent\textit{Source reference.} Original reference 07DW: focused-seeds/geometry/S9-smo-resolve-general/ref-tags.txt, source lines 8--8. Navigation pointer only: the counted proof cites this statement by \textbackslash{}ref as a standing input; the deduction the counted proof has to make is carried out inside the counted span, and no part of that deduction is delegated to this pointer.\par\medskip
\ifdefined\hypertarget\hypertarget{sourcepointer-5}{}\fi
\noindent\textit{Source reference.} Original reference 07DV: focused-seeds/geometry/S9-smo-resolve-general/ref-tags.txt, source lines 9--9. Navigation pointer only: the counted proof cites this statement by \textbackslash{}ref as a standing input; the deduction the counted proof has to make is carried out inside the counted span, and no part of that deduction is delegated to this pointer.\par\medskip
\ifdefined\hypertarget\hypertarget{sourcepointer-6}{}\fi
\noindent\textit{Source reference.} Original lemma 07EZ: sources/stacks-smoothing/pinned-original-source.tex, source lines 936--936. Navigation pointer only: the counted proof cites this statement by \textbackslash{}ref as a standing input; the deduction the counted proof has to make is carried out inside the counted span, and no part of that deduction is delegated to this pointer.\par\medskip
\ifdefined\hypertarget\hypertarget{sourcepointer-7}{}\fi
\noindent\textit{Source reference.} Original lemma 07FA: sources/stacks-smoothing/pinned-original-source.tex, source lines 2447--2447. Navigation pointer only: the counted proof cites this statement by \textbackslash{}ref as a standing input; the deduction the counted proof has to make is carried out inside the counted span, and no part of that deduction is delegated to this pointer.\par\medskip
\ifdefined\hypertarget\hypertarget{sourcepointer-8}{}\fi
\noindent\textit{Source reference.} Original lemma 07FI: sources/stacks-smoothing/pinned-original-source.tex, source lines 2841--2841. Navigation pointer only: the counted proof cites this statement by \textbackslash{}ref as a standing input; the deduction the counted proof has to make is carried out inside the counted span, and no part of that deduction is delegated to this pointer.\par\medskip
\ifdefined\hypertarget\hypertarget{sourcepointer-9}{}\fi
\noindent\textit{Source reference.} Original lemma 07CE: sources/stacks-smoothing/pinned-original-source.tex, source lines 559--559. Navigation pointer only: the counted proof cites this statement by \textbackslash{}ref as a standing input; the deduction the counted proof has to make is carried out inside the counted span, and no part of that deduction is delegated to this pointer.\par\medskip
\ifdefined\hypertarget\hypertarget{sourcepointer-10}{}\fi
\noindent\textit{Source reference.} Original lemma 07F8: sources/stacks-smoothing/pinned-original-source.tex, source lines 2314--2314. Navigation pointer only: the counted proof cites this statement by \textbackslash{}ref as a standing input; the deduction the counted proof has to make is carried out inside the counted span, and no part of that deduction is delegated to this pointer.\par\medskip
\ifdefined\hypertarget\hypertarget{sourcepointer-11}{}\fi
\noindent\textit{Source reference.} Original lemma 07FC: sources/stacks-smoothing/pinned-original-source.tex, source lines 2504--2504. Navigation pointer only: the counted proof cites this statement by \textbackslash{}ref as a standing input; the deduction the counted proof has to make is carried out inside the counted span, and no part of that deduction is delegated to this pointer.\par\medskip
\ifdefined\hypertarget\hypertarget{sourcepointer-12}{}\fi
\noindent\textit{Source reference.} Original lemma 07FH: sources/stacks-smoothing/pinned-original-source.tex, source lines 2769--2769. Navigation pointer only: the counted proof cites this statement by \textbackslash{}ref as a standing input; the deduction the counted proof has to make is carried out inside the counted span, and no part of that deduction is delegated to this pointer.\par\medskip
\ifdefined\hypertarget\hypertarget{sourcepointer-13}{}\fi
\noindent\textit{Source reference.} Original lemma 07CC: sources/stacks-smoothing/pinned-original-source.tex, source lines 512--512. Navigation pointer only: the counted proof cites this statement by \textbackslash{}ref as a standing input; the deduction the counted proof has to make is carried out inside the counted span, and no part of that deduction is delegated to this pointer.\par\medskip
\ifdefined\hypertarget\hypertarget{sourcepointer-14}{}\fi
\noindent\textit{Source reference.} Original reference 07F7: sources/stacks-smoothing/pinned-original-source.tex, source lines 2296--2296. Navigation pointer only: the counted proof cites this statement by \textbackslash{}ref as a standing input; the deduction the counted proof has to make is carried out inside the counted span, and no part of that deduction is delegated to this pointer.\par\medskip
\begin{lemma}
\label{lemma-resolve-general}
Let $k \to A \to \Lambda \supset \mathfrak q$ be as in
Situation \ref{situation-local} where
\begin{enumerate}
\item $k$ is a field of characteristic $p > 0$,
\item $\Lambda$ is Noetherian and geometrically regular over $k$,
\item $\mathfrak q$ is minimal over $\mathfrak h_A$.
\end{enumerate}
Then $k \to A \to \Lambda \supset \mathfrak q$ can be resolved.
\end{lemma}

\begin{proof}
The lemma is proven by the following steps in the given order.
We will justify each of these steps below.
\begin{enumerate}
\item
\label{item-power}
Pick an integer $N > 0$ such that
$\mathfrak q^N\Lambda_\mathfrak q \subset H_{A/k}\Lambda_\mathfrak q$.
\item
\label{item-generators}
Pick generators $a_1, \ldots, a_t \in A$ of the ideal $H_{A/R}$.
\item
\label{item-dimension}
Set $d = \dim(\Lambda_\mathfrak q)$.
\item
\label{item-standardizer}
Set $B = A[x_1, \ldots, x_d, z_{ij}]/(x_i^{2N} - \sum z_{ij}a_j)$.
\item
\label{item-elkik}
Consider $B$ as a $k[x_1, \ldots, x_d]$-algebra and let
$B \to C$ be as in
Lemma \ref{lemma-improve-presentation}.
We also obtain a section $C \to B$.
\item
\label{item-strictly-standard}
Choose $c > 0$ such that each $x_i^c$
is strictly standard in $C$ over $k[x_1, \ldots, x_d]$.
\item
\label{item-set-n}
Set $n = N + dc$ and $e = 8c$.
\item
\label{item-set-E}
Let $E \subset \Lambda_\mathfrak q/\mathfrak q^n\Lambda_\mathfrak q$
be the images of generators of $A$ as a $k$-algebra.
\item
\label{item-NP}
Choose an integer $m$ and a $k$-algebra map
$\varphi : k[y_1, \ldots, y_m] \to \Lambda$
and a factorization by local Artinian rings
$$
k[y_1, \ldots, y_m]_\mathfrak p/\mathfrak p^n k[y_1, \ldots, y_m]_\mathfrak p
\to D \to
\Lambda_\mathfrak q/\mathfrak q^n\Lambda_\mathfrak q
$$
such that the first arrow is essentially smooth, the second is flat,
$E$ is contained in $D$, with $\mathfrak p = \varphi^{-1}(\mathfrak q)$
the map $k[y_1, \ldots, y_m]_\mathfrak p \to \Lambda_\mathfrak q$ is
flat, and $\mathfrak p \Lambda_\mathfrak q = \mathfrak q \Lambda_\mathfrak q$.
\item
\label{item-choose-pii}
Choose $\pi_1, \ldots, \pi_d \in \mathfrak p$ which map to a
regular system of parameters of $k[y_1, \ldots, y_m]_\mathfrak p$.
\item
\label{item-set-R}
Let $R = k[y_1, \ldots, y_m, t_1, \ldots, t_m]$ and $\gamma_i = \pi_i t_i$.
\item
\label{item-modify-pii}
If necessary modify the choice of $\pi_i$ such that
for $i = 1, \ldots, d$ we have 
$$
\text{Ann}_{R/(\gamma_1^e, \ldots, \gamma_{i - 1}^e)R}(\gamma_i)
=
\text{Ann}_{R/(\gamma_1^e, \ldots, \gamma_{i - 1}^e)R}(\gamma_i^2)
$$
\item
\label{item-choose-deltai}
There exist $\delta_1, \ldots, \delta_d \in \Lambda$,
$\delta_i \not \in \mathfrak q$ and a factorization
$D \to D' \to \Lambda_\mathfrak q/\mathfrak q^n\Lambda_\mathfrak q$
with $D'$ local Artinian, $D \to D'$ essentially smooth, the map
$D' \to \Lambda_\mathfrak q/\mathfrak q^n\Lambda_\mathfrak q$
flat such that, with $\pi_i' = \delta_i \pi_i$, we have for
$i = 1, \ldots, d$
\begin{enumerate}
\item $(\pi_i')^{2N} = \sum a_j\lambda_{ij}$ in $\Lambda$ where
$\lambda_{ij} \bmod \mathfrak q^n\Lambda_\mathfrak q$ is an element of $D'$,
\item $\text{Ann}_{\Lambda/({\pi'}_1^e, \ldots, {\pi'}_{i - 1}^e)}({\pi'}_i) =
\text{Ann}_{\Lambda/({\pi'}_1^e, \ldots, {\pi'}_{i - 1}^e)}({\pi'}_i^2)$,
\item $\delta_i \bmod \mathfrak q^n\Lambda_\mathfrak q$ is an element of $D'$.
\end{enumerate}
\item
\label{item-map-B-Lambda}
Define $B \to \Lambda$ by sending $x_i$ to $\pi'_i$ and
$z_{ij}$ to $\lambda_{ij}$ found above. Define $C \to \Lambda$
by composing the map $B \to \Lambda$ with the retraction $C \to B$.
\item
\label{item-set-map-R}
Map $R \to \Lambda$ by $\varphi$ on $k[y_1, \ldots, y_m]$
and by sending $t_i$ to $\delta_i$. Further introduce a map
$$
k[x_1, \ldots, x_d]
\longrightarrow
R = k[y_1, \ldots, y_m, t_1, \ldots, t_d]
$$
by sending $x_i$ to $\gamma_i = \pi_i t_i$.
\item
\label{item-first-solve}
It suffices to resolve
$$
R
\to
C \otimes_{k[x_1, \ldots, x_d]} R
\to
\Lambda \supset \mathfrak q
$$
\item
\label{item-set-I}
Set $I = (\gamma_1^e, \ldots, \gamma_d^e) \subset R$.
\item
\label{item-second-resolve}
It suffices to resolve
$$
R/I
\to
C \otimes_{k[x_1, \ldots, x_d]} R/I
\to
\Lambda/I\Lambda \supset \mathfrak q/I\Lambda
$$
\item
\label{item-set-r}
We denote $\mathfrak r \subset R = k[y_1, \ldots, y_m, t_1, \ldots, t_d]$
the inverse image of $\mathfrak q$.
\item
\label{item-third-resolve}
It suffices to resolve
$$
(R/I)_\mathfrak r \to
C \otimes_{k[x_1, \ldots, x_d]} (R/I)_\mathfrak r \to
\Lambda_\mathfrak q/I\Lambda_\mathfrak q
\supset
\mathfrak q\Lambda_\mathfrak q/I\Lambda_\mathfrak q
$$
\item
\label{item-fourth-resolve}
Set $J = (\pi_1^e, \ldots, \pi_d^e)$ in $k[y_1, \ldots, y_m]$.
\item
\label{item-fifth-resolve}
It suffices to resolve
$$
(R/JR)_\mathfrak p \to
C \otimes_{k[x_1, \ldots, x_d]} (R/JR)_\mathfrak p \to
\Lambda_\mathfrak q/J\Lambda_\mathfrak q
\supset
\mathfrak q\Lambda_\mathfrak q/J\Lambda_\mathfrak q
$$
\item
\label{item-sixth-resolve}
It suffices to resolve
$$
(R/\mathfrak p^nR)_\mathfrak p \to
C \otimes_{k[x_1, \ldots, x_d]} (R/\mathfrak p^nR)_\mathfrak p \to
\Lambda_\mathfrak q/\mathfrak q^n\Lambda_\mathfrak q
\supset
\mathfrak q\Lambda_\mathfrak q/\mathfrak q^n\Lambda_\mathfrak q
$$
\item
\label{item-seventh-resolve}
It suffices to resolve
$$
(R/\mathfrak p^nR)_\mathfrak p \to
B \otimes_{k[x_1, \ldots, x_d]} (R/\mathfrak p^nR)_\mathfrak p \to
\Lambda_\mathfrak q/\mathfrak q^n\Lambda_\mathfrak q
\supset
\mathfrak q\Lambda_\mathfrak q/\mathfrak q^n\Lambda_\mathfrak q
$$
\item
\label{item-eighth-resolve}
The ring $D'[t_1, \ldots, t_d]$ is given the structure of an
$R_\mathfrak p/\mathfrak p^nR_\mathfrak p$-algebra by the given map
$k[y_1, \ldots, y_m]_\mathfrak p/\mathfrak p^n k[y_1, \ldots, y_m]_\mathfrak p
\to D'$ and by sending $t_i$ to $t_i$. It suffices to find a factorization
$$
B \otimes_{k[x_1, \ldots, x_d]} (R/\mathfrak p^nR)_\mathfrak p
\to D'[t_1, \ldots, t_d] \to
\Lambda_\mathfrak q/\mathfrak q^n\Lambda_\mathfrak q
$$
where the second arrow sends $t_i$ to $\delta_i$ and induces the given
homomorphism $D' \to \Lambda_\mathfrak q/\mathfrak q^n\Lambda_\mathfrak q$.
\item
\label{item-done}
Such a factorization exists by our choice of $D'$ above.
\end{enumerate}
We now give the justification for each of the steps, except that we
skip justifying the steps which just introduce notation.

\medskip\noindent
Ad (\ref{item-power}). This is possible as $\mathfrak q$ is minimal
over $\mathfrak h_A = \sqrt{H_{A/k}\Lambda}$.

\medskip\noindent
Ad (\ref{item-strictly-standard}). Note that $A_{a_i}$ is smooth
over $k$. Hence $B_{a_j}$, which is isomorphic to a polynomial
algebra over $A_{a_j}[x_1, \ldots, x_d]$, is smooth over
$k[x_1, \ldots, x_d]$. Thus $B_{x_i}$ is smooth over $k[x_1, \ldots, x_d]$.
By Lemma \ref{lemma-improve-presentation}
we see that $C_{x_i}$ is smooth over $k[x_1, \ldots, x_d]$
with finite free module of differentials. Hence some power of
$x_i$ is strictly standard in $C$ over $k[x_1, \ldots, x_n]$
by Lemma \ref{lemma-compare-standard}.

\medskip\noindent
Ad (\ref{item-NP}). This follows by applying Lemma \ref{lemma-solution-modulo}.

\medskip\noindent
Ad (\ref{item-choose-pii}). Since
$k[y_1, \ldots, y_m]_\mathfrak p \to \Lambda_\mathfrak q$ is
flat and $\mathfrak p \Lambda_\mathfrak q = \mathfrak q \Lambda_\mathfrak q$
by construction
we see that $\dim(k[y_1, \ldots, y_m]_\mathfrak p) = d$ by
Algebra, Lemma \ref{algebra-lemma-dimension-base-fibre-equals-total}.
Thus we can find $\pi_1, \ldots, \pi_d \in \Lambda$ which map to
a regular system of parameters in $\Lambda_\mathfrak q$.

\medskip\noindent
Ad (\ref{item-modify-pii}). By
Algebra, Lemma \ref{algebra-lemma-regular-ring-CM}
any permutation of the sequence $\pi_1, \ldots, \pi_d$ is a
regular sequence in $k[y_1, \ldots, y_m]_\mathfrak p$. Hence
$\gamma_1 = \pi_1 t_1, \ldots, \gamma_d = \pi_d t_d$ is a regular
sequence in
$R_\mathfrak p = k[y_1, \ldots, y_m]_\mathfrak p[t_1, \ldots, t_d]$, see
Algebra, Lemma \ref{algebra-lemma-regular-sequence-in-polynomial-ring}.
Let $S = k[y_1, \ldots, y_m] \setminus \mathfrak p$ so that
$R_\mathfrak p = S^{-1}R$. Note that $\pi_1, \ldots, \pi_d$
and $\gamma_1, \ldots, \gamma_d$
remain regular sequences if we multiply our $\pi_i$ by elements of $S$.
Suppose that
$$
\text{Ann}_{R/(\gamma_1^e, \ldots, \gamma_{i - 1}^e)R}(\gamma_i)
=
\text{Ann}_{R/(\gamma_1^e, \ldots, \gamma_{i - 1}^e)R}(\gamma_i^2)
$$
holds for $i = 1, \ldots, t$ for some $t \in \{0, \ldots, d\}$. Note that
$\gamma_1^e, \ldots, \gamma_t^e, \gamma_{t + 1}$ is a
regular sequence in $S^{-1}R$ by
Algebra, Lemma \ref{algebra-lemma-regular-sequence-powers}.
Hence we see that
$$
\text{Ann}_{S^{-1}R/(\gamma_1^e, \ldots, \gamma_{i - 1}^e)}(\gamma_i) =
\text{Ann}_{S^{-1}R/(\gamma_1^e, \ldots, \gamma_{i - 1}^e)}(\gamma_i^2).
$$
Thus we get
$$
\text{Ann}_{R/(\gamma_1^e, \ldots, \gamma_t^e)R}(\gamma_{t + 1})
=
\text{Ann}_{R/(\gamma_1^e, \ldots, \gamma_t^e)R}(\gamma_{t + 1}^2)
$$
after replacing $\pi_{t + 1}$ by $s\pi_{t + 1}$ for some $s \in S$ by
Lemma \ref{lemma-ogoma}. By induction on $t$ this produces the desired
sequence.

\medskip\noindent
Ad (\ref{item-choose-deltai}). Let $S = \Lambda \setminus \mathfrak q$
so that $\Lambda_\mathfrak q = S^{-1}\Lambda$. Set
$\bar \Lambda = \Lambda_\mathfrak q/\mathfrak q^n \Lambda_\mathfrak q$.
Suppose that we have a $t \in \{0, \ldots, d\}$ and
$\delta_1, \ldots, \delta_t \in S$ and a factorization
$D \to D' \to \bar \Lambda$ as in (\ref{item-choose-deltai})
such that (a), (b), (c) hold for $i = 1, \ldots, t$. We have
$\pi_{t + 1}^N \in H_{A/k}\Lambda_\mathfrak q$
as $\mathfrak q^N \Lambda_\mathfrak q \subset H_{A/k}\Lambda_\mathfrak q$
by (\ref{item-power}). Hence
$\pi_{t + 1}^N \in H_{A/k} \bar\Lambda$. Hence
$\pi_{t + 1}^N \in H_{A/k}D'$ as $D' \to \bar \Lambda$
is faithfully flat, see
Algebra, Lemma \ref{algebra-lemma-faithfully-flat-universally-injective}.
Recall that $H_{A/k} = (a_1, \ldots, a_t)$.
Say $\pi_{t + 1}^N = \sum a_j d_j$ in $D'$ and choose
$c_j \in \Lambda_\mathfrak q$ lifting $d_j \in D'$. Then
$\pi_{t + 1}^N = \sum c_j a_j + \epsilon$ with
$\epsilon \in \mathfrak q^n\Lambda_\mathfrak q \subset
\mathfrak q^{n - N}H_{A/k}\Lambda_\mathfrak q$.
Write $\epsilon = \sum a_j c'_j$ for some
$c'_j \in \mathfrak q^{n - N}\Lambda_\mathfrak q$.
Hence $\pi_{t + 1}^{2N} = \sum (\pi_{t + 1}^N c_j + \pi_{t + 1}^N c'_j) a_j$.
Note that $\pi_{t + 1}^Nc'_j$ maps to zero in $\bar \Lambda$; this trivial
but key observation will ensure later that (a) holds.
Now we choose $s \in S$ such that there exist
$\mu_{t + 1j} \in \Lambda$ such that on the one hand
$\pi_{t + 1}^N c_j + \pi_{t + 1}^N c'_j = \mu_{t + 1j}/s^{2N}$
in $S^{-1}\Lambda$ and on the other
$(s \pi_{t + 1})^{2N} = \sum \mu_{t + 1j}a_j$
in $\Lambda$ (minor detail omitted). We may further replace $s$ by
a power and enlarge $D'$ such that $s$ maps to an element of $D'$.
With these choices $\mu_{t + 1j}$ maps to $s^{2N}d_j$ which is
an element of $D'$. Note that $\pi_1, \ldots, \pi_d$ are a regular
sequence of parameters in $S^{-1}\Lambda$ by our
choice of $\varphi$. Hence $\pi_1, \ldots, \pi_d$ forms a regular sequence
in $\Lambda_\mathfrak q$ by
Algebra, Lemma \ref{algebra-lemma-regular-ring-CM}.
It follows that ${\pi'}_1^e, \ldots, {\pi'}_t^e, s\pi_{t + 1}$ is a
regular sequence in $S^{-1}\Lambda$ by
Algebra, Lemma \ref{algebra-lemma-regular-sequence-powers}.
Thus we get
$$
\text{Ann}_{S^{-1}\Lambda/({\pi'}_1^e, \ldots, {\pi'}_t^e)}(s\pi_{t + 1}) =
\text{Ann}_{S^{-1}\Lambda/({\pi'}_1^e, \ldots, {\pi'}_t^e)}((s\pi_{t + 1})^2).
$$
Hence we may apply Lemma \ref{lemma-ogoma} to find an $s' \in S$
such that
$$
\text{Ann}_{\Lambda/({\pi'}_1^e, \ldots, {\pi'}_t^e)}((s')^qs\pi_{t + 1})
=
\text{Ann}_{\Lambda/({\pi'}_1^e, \ldots, {\pi'}_t^e)}(((s')^qs\pi_{t + 1})^2).
$$
for any $q > 0$. By Lemma \ref{lemma-enlarge-solution-modulo}
we can choose $q$ and enlarge $D'$ such that $(s')^q$ maps to an element
of $D'$. Setting $\delta_{t + 1} = (s')^qs$ and we conclude that
(a), (b), (c) hold for $i = 1, \ldots, t + 1$. For (a) note that
$\lambda_{t + 1j} = (s')^{2Nq}\mu_{t + 1j}$ works.
By induction on $t$ we win.

\medskip\noindent
Ad (\ref{item-first-solve}). By construction the radical of
$H_{(C \otimes_{k[x_1, \ldots, x_d]} R)/R} \Lambda$ contains
$\mathfrak h_A$. Namely, the elements $a_j \in H_{A/k}$
map to elements of $H_{B/k[x_1, \ldots, x_n]}$, hence map to elements
of $H_{C/k[x_1, \ldots, x_n]}$, hence $a_j \otimes 1$ map to elements of
$H_{C \otimes_{k[x_1, \ldots, x_d]} R/R}$. Moreover, if we have a solution
$C \otimes_{k[x_1, \ldots, x_n]} R \to T \to \Lambda$ of
$$
R
\to
C \otimes_{k[x_1, \ldots, x_d]} R
\to
\Lambda \supset \mathfrak q
$$
then $H_{T/R} \subset H_{T/k}$ as $R$ is smooth over $k$.
Hence $T$ will also be a solution for
the original situation $k \to A \to \Lambda \supset \mathfrak q$.

\medskip\noindent
Ad (\ref{item-second-resolve}). Follows on applying
Lemma \ref{lemma-lift-solution} to
$R \to C \otimes_{k[x_1, \ldots, x_d]} R
\to \Lambda \supset \mathfrak q$ and the sequence of
elements $\gamma_1^c, \ldots, \gamma_d^c$. We note that since $x_i^c$
are strictly standard in $C$ over $k[x_1, \ldots, x_d]$ the elements
$\gamma_i^c$ are strictly standard in $C \otimes_{k[x_1, \ldots, x_d]} R$
over $R$ by Lemma \ref{lemma-strictly-standard-base-change}.
The other assumption of Lemma \ref{lemma-lift-solution} holds by steps
(\ref{item-modify-pii}) and (\ref{item-choose-deltai}).

\medskip\noindent
Ad (\ref{item-third-resolve}). Apply Lemma \ref{lemma-delocalize-height-zero}
to the situation in (\ref{item-second-resolve}). In the rest of the
arguments the target ring is local Artinian, hence we are looking for
a factorization by a smooth algebra $T$ over the source ring.

\medskip\noindent
Ad (\ref{item-fifth-resolve}).
Suppose that $C \otimes_{k[x_1, \ldots, x_d]} (R/JR)_\mathfrak p \to
T \to \Lambda_\mathfrak q/J\Lambda_\mathfrak q$ is a solution to
$$
(R/JR)_\mathfrak p \to
C \otimes_{k[x_1, \ldots, x_d]} (R/JR)_\mathfrak p \to
\Lambda_\mathfrak q/J\Lambda_\mathfrak q
\supset
\mathfrak q\Lambda_\mathfrak q/J\Lambda_\mathfrak q
$$
Then $C \otimes_{k[x_1, \ldots, x_d]} (R/I)_\mathfrak r \to T_\mathfrak r \to
\Lambda_\mathfrak q/I\Lambda_\mathfrak q$
is a solution to the situation in (\ref{item-third-resolve}).

\medskip\noindent
Ad (\ref{item-sixth-resolve}). Our $n = N + dc$ is large enough so that
$\mathfrak p^nk[y_1, \ldots, y_m]_\mathfrak p \subset J_\mathfrak p$
and $\mathfrak q^n \Lambda_\mathfrak q \subset J\Lambda_\mathfrak q$.
Hence if we have a solution
$C \otimes_{k[x_1, \ldots, x_d]} (R/\mathfrak p^nR)_\mathfrak p \to
T \to \Lambda_\mathfrak q/\mathfrak q^n\Lambda_\mathfrak q$
of (\ref{item-fifth-resolve}
then we can take $T/JT$ as the solution for
(\ref{item-sixth-resolve}).

\medskip\noindent
Ad (\ref{item-seventh-resolve}). This is true because we have a
section $C \to B$ in the category of $R$-algebras.

\medskip\noindent
Ad (\ref{item-eighth-resolve}). This is true because $D'$ is
essentially smooth over the local Artinian ring
$k[y_1, \ldots, y_m]_\mathfrak p/\mathfrak p^n k[y_1, \ldots, y_m]_\mathfrak p$
and
$$
R_\mathfrak p/\mathfrak p^nR_\mathfrak p =
k[y_1, \ldots, y_m]_\mathfrak p/
\mathfrak p^n k[y_1, \ldots, y_m]_\mathfrak p[t_1, \ldots, t_d].
$$
Hence $D'[t_1, \ldots, t_d]$ is a filtered colimit of smooth
$R_\mathfrak p/\mathfrak p^nR_\mathfrak p$-algebras and
$B \otimes_{k[x_1, \ldots, x_d]} (R_\mathfrak p/\mathfrak p^nR_\mathfrak p)$
factors through one of these.

\medskip\noindent
Ad (\ref{item-done}). The final twist of the proof is that we cannot
just use the map $B \to D'$ which maps $x_i$ to the image of $\pi_i'$
in $D'$ and $z_{ij}$ to the image of $\lambda_{ij}$ in $D'$
because we need the diagram
$$
\xymatrix{
B \ar[r] & D'[t_1, \ldots, t_d] \\
k[x_1, \ldots, x_d] \ar[r] \ar[u] &
R_\mathfrak p/\mathfrak p^nR_\mathfrak p \ar[u]
}
$$
to commute and we need the composition
$B \to D'[t_1, \ldots, t_d] \to
\Lambda_\mathfrak q/\mathfrak q^n\Lambda_\mathfrak q$
to be the map of (\ref{item-map-B-Lambda}).
This requires us to map $x_i$ to the image of
$\pi_i t_i$ in $D'[t_1, \ldots, t_d]$.
Hence we map $z_{ij}$ to the image of
$\lambda_{ij} t_i^{2N} / \delta_i^{2N}$ in $D'[t_1, \ldots, t_d]$
and everything is clear.
\end{proof}

\end{document}
